\section{Notaci\'on y convenciones}
\label{sec:notacion}

Estos apuntes cubren dos bloques: c\'alculo vectorial y geometr\'ia diferencial
(Secciones \ref{sec:espacios_vectoriales}--\ref{sec:formas_diferenciales}) y
variable compleja, series y an\'alisis de Fourier
(Secciones \ref{sec:variable_compleja}--\ref{sec:integracion_compleja}).

\begin{itemize}
  \item $V$ denota un espacio vectorial sobre el campo $\F$, con
        $\F=\R$ o $\F=\C$; $\Ee^n$ denota el espacio eucl\'ideo de
        dimensi\'on $n$.
  \item $\vb{a},\vb{b},\vb{c}$ son vectores y $\vh{e}_i$ los elementos de una
        base ortonormal. La transposici\'on se escribe $\vb{a}^\tp$ y la
        transpuesta conjugada $\vb{a}^\dg$.
  \item $\wedge$ es el \textbf{producto cu\~na} o exterior; en $\R^3$ coincide
        con el producto cruz. El producto punto se escribe $\vb{a}\cdot\vb{b}$
        y el producto interno abstracto $g(\vb{a},\vb{b})$ o
        $\braket{\vb{a}}{\vb{b}}$.
  \item $\otimes$ es el producto directo y $\oplus$ la suma directa.
  \item $\vb{r}(t)$ denota una curva y $s$ su longitud de arco; la prima
        $\vb{T}'$ siempre significa derivada respecto de $s$.
  \item $\vb{M}(\xi,\et)$ denota una superficie parametrizada y
        $\vb{M}(\xi,\et,\zeta)$ un mapeo de volumen.
  \item En la parte compleja, $z=x+iy$, $\bar z$ es el conjugado,
        $\abs{z}=(z\bar z)^{1/2}$ y $\vph=\Argu(z)$ el argumento principal.
        $\Om$ y $D$ denotan dominios abiertos de $\C$ y $\ga$ una curva
        parametrizada contenida en ellos.
  \item $\lito(\vb{u})$ es la $\lito$ peque\~na de Landau, definida por
        $\lim_{\vb{u}\rightarrow\vb{0}}\norm{\lito(\vb{u})}/\norm{\vb{u}}=0$.
\end{itemize}

\subsection{Notaci\'on de Einstein}
\label{subsec:einstein}

Un \'indice repetido dentro de un mismo t\'ermino indica suma sobre todo su
rango:
\begin{equation}
  \label{eq:einstein}
  \sum_{i=1}^{n}a_ib_i=a_1b_1+a_2b_2+\dots+a_nb_n=a_ib_i .
\end{equation}
As\'i, la serie de Taylor $f(x)=\sum_{n=0}^{\infty}a_n(x-x_0)^n$ se abrevia
$f(x)=a_n(x-x_0)^n$.

\subsection{Delta de Kronecker}
\label{subsec:kronecker}

Sea $\set{\vh{e}_i}_{i=1}^{n}$ una base ortonormal de $\Ee^n$, es decir,
$\vh{e}_k$ tiene un $1$ en la entrada $k$-\'esima y ceros en las dem\'as.

\begin{definicion}[Delta de Kronecker]
  \label{def:kronecker}
  \begin{equation}
    \label{eq:kronecker}
    \de_{ij}=g(\vh{e}_i,\vh{e}_j)=\pdv{x_j}{x_i}=
    \begin{cases}
      0, & i\neq j,\\
      1, & i=j .
    \end{cases}
  \end{equation}
\end{definicion}

Su utilidad es \emph{contraer} sumas. Si $\vb{a}=\sum_{i=1}^{n}a_i\vh{e}_i$,
entonces
\begin{equation}
  \label{eq:contraccion_kronecker}
  \begin{aligned}
    g(\vb{a},\vh{e}_j)
      &=\qty{\sum_{i=1}^{n}a_i\vh{e}_i}\cdot\vh{e}_j
       =\sum_{i=1}^{n}a_i\,\vh{e}_i\cdot\vh{e}_j\\
      &=\sum_{i=1}^{n}a_i\de_{ij}=a_j ,
  \end{aligned}
\end{equation}
ya que todos los sumandos con $i\neq j$ se anulan.

\subsection{S\'imbolos de Levi-Civita}
\label{subsec:levi_civita}

Tres vectores linealmente independientes generan una caja de volumen
\begin{equation}
  \label{eq:volumen_caja}
  \mathrm{Vol}=(\text{\'area de la base})(\text{altura})
    =(\vb{a}\wedge\vb{b})\cdot\vb{c},
\end{equation}
cuyo signo define la \textbf{orientaci\'on} del sistema: es positiva si
$(\vb{a}\wedge\vb{b})\cdot\vb{c}>0$ y negativa en caso contrario.

\begin{definicion}[S\'imbolo de Levi-Civita]
  \label{def:levi_civita}
  Sea $\set{\vh{e}_i}_{i=1}^{3}$ la base can\'onica de $\Ee^3$. Entonces
  \begin{equation}
    \label{eq:levi_civita}
    \ep_{ijk}=(\vh{e}_i\wedge\vh{e}_j)\cdot\vh{e}_k=
    \begin{cases}
      0, & \text{si se repite alg\'un \'indice},\\
      -1, & \text{si la permutaci\'on es impar},\\
      1, & \text{si la permutaci\'on es par}.
    \end{cases}
  \end{equation}
\end{definicion}

Con \'el, la componente $i$-\'esima del producto cu\~na en $\R^3$ se escribe en
una sola l\'inea:
\begin{equation}
  \label{eq:cuna_levicivita}
  (\vb{a}\wedge\vb{b})_i=\ep_{ijk}a_jb_k .
\end{equation}

\begin{ejemplo}
  \label{ej:cuna_componentes}
  Para $i=1$ los \'unicos t\'erminos no nulos son $\ep_{123}$ y $\ep_{132}$:
  \begin{equation}
    \label{eq:cuna_primera_componente}
    \begin{aligned}
      (\vb{a}\wedge\vb{b})_1
        &=\ep_{123}a_2b_3+\ep_{132}a_3b_2\\
        &=a_2b_3-a_3b_2 ,
    \end{aligned}
  \end{equation}
  y an\'alogamente para las otras dos, de modo que
  \begin{equation}
    \label{eq:cuna_3d}
    \vb{a}\wedge\vb{b}=
    \begin{pmatrix}
      a_2b_3-a_3b_2\\
      a_3b_1-a_1b_3\\
      a_1b_2-a_2b_1
    \end{pmatrix}.
  \end{equation}
\end{ejemplo}

El triple producto y el rotacional admiten la misma abreviatura:
\begin{equation}
  \label{eq:triple_producto}
  (\vb{a}\wedge\vb{b})\cdot\vb{c}=\ep_{ijk}c_ia_jb_k,
  \qquad
  (\rot\vb{f})_i=\ep_{ijk}\partial_jf_k .
\end{equation}
